% Formal checks for the explicit editorial reconstructions, not an Aristotle calculus.
\section{Refutation and inferential diagnosis: countermodels and proofs}
\label{sec:annex-p18-refutation-diagnosis}

The interpretations below use
Convention~\ref{def:p18-rd-formal-conventions}. They verify the
chapter's explicit modern reconstructions. The historical categories
and responses remain attributed to the selected translation
\cite{AristotleSophisticalRefutationsPickardCambridgeMIT}.

\begin{annexitem}{annex:p18.rd.refutation-standard}{thm:p18-rd-refutation-standard}{Refutation, invalidity and truth}
\begin{proof}
Suppose \(\Gamma\models_{\mathcal K}D\) and
\(D\models_{\mathcal K}\neg C\). Any interpretation in
\(\mathcal K\) satisfying \(\Gamma\) satisfies \(D\) by
the first premise, and hence \(\neg C\) by the second. This is
exactly \(\Gamma\models_{\mathcal K}\neg C\). Applying this
conditional result to an actual case requires an interpretation in
the specified class in which the premises are true. The consequence
relation alone does not assert that the actual case has these properties.

For failure of the inference from invalidity to falsity, take distinct
propositional letters \(p,q\), premises \(\Gamma=\{p\}\)
and conclusion \(C=q\), over all Boolean valuations. The valuation
\(p=1,q=0\) satisfies the premise and falsifies the conclusion, so
\(\Gamma\not\models q\). The valuation \(p=1,q=1\) satisfies
the same premise and falsifies \(\neg q\), so
\(\Gamma\not\models\neg q\) as well. In the second valuation
the invalid argument has both a true premise and a true conclusion.
Validity concerns all admitted interpretations; truth at one actual
interpretation is a different property.

An independently established \(\neg C\) concerns the target's
truth. It does not identify which step or premise in a particular
argument for \(C\) failed. Different invalid arguments can have
the same conclusion while failing for different reasons. The explicit
countermodels in the following proofs locate the defect in each
specified reconstruction rather than relying only on a contrary
conclusion.

The warning-source example is the same distinction with meaningful
propositional labels. Take \(\Gamma=\{u\}\) and \(C=w\).
Both assignments \((u,w)=(1,1)\) and \((1,0)\) satisfy the
source-unreliability premise. One makes the warning true and the
other false. Thus that premise entails neither \(w\) nor
\(\neg w\), without an additional connection to the report's
actual subject matter.
\end{proof}
\end{annexitem}

\begin{annexitem}{annex:p18.rd.language-cases}{thm:p18-rd-language-cases}{Explicit models for the language diagnoses}
\begin{proof}
\emph{Changed middle term.} Use a domain of two distinct objects
\(a,b\), interpret \(F\) by \(\{a\}\), \(R\) by
\(\{b\}\), and \(W\) by \(\{b\}\). Then \(F(a)\)
is true. Every river-bank object is in \(W\), so
\(\forall x(R(x)\to W(x))\) is true. But \(W(a)\) is false.
This countermodel does not rely on an empty river-bank class. If a
further premise \(R(a)\) holds in any interpretation of the
universal premise, its instance at \(a\) forces \(W(a)\).
This proves the validity of the conditional repair, not the truth
of its additional factual premise.

\emph{Sight and argument place.} Take domain \(\{o,t\}\)
with \(o\ne t\), relation \(S=\{(o,t)\}\) and
\(V=\{o\}\). The observed relation \(S(o,t)\) is true.
The only true antecedent of an instance of
\(S(x,y)\to V(x)\) is at \((o,t)\), and its consequent
is true, so the universal premise holds. Yet \(V(t)\) is false.
In any interpretation satisfying those premises, the universal
instance at \(o,t\) and the truth of \(S(o,t)\) force
\(V(o)\). The two proposed conclusions thus have different
statuses under the same fixed interpretation of seeing.

\emph{Ability scope.} Let the worlds be \(w_0,w_1\), with
the sole accessibility pair \((w_0,w_1)\). Make \(S\) true
only at \(w_0\), and \(W\) true only at \(w_1\).
At \(w_0\), \(S\) is true and the accessible world \(w_1\)
makes \(W\) true, so \(S\land\Diamond W\) holds.
The sole accessible world makes \(S\) false, so no accessible
world makes \(S\land W\) true. Hence
\(\Diamond(S\land W)\) is false at \(w_0\).
The retained conclusion \(\Diamond W\) follows in every model
from the truth condition of conjunction.

\emph{Odd sum.} In the natural numbers \(5=2+3\), and
\(5=2\cdot2+1\), so both premises are true. There is no natural
\(k\) with \(2=2k+1\): \(k=0\) gives one, and
\(k\geq1\) gives at least three. Thus \(\mathrm{Odd}(2)\)
is false and the proposed conjunction of odd summands is false.
Replacing the equal argument \(5\) by \(2+3\) in the
original predicate gives the true retained conclusion
\(\mathrm{Odd}(2+3)\). Equality substitution does not distribute
the predicate over addition.

\emph{Lost distinction in transcription.} Let the token set be
\(\{e_1,e_2\}\) with distinct members. Let \(f\) map both
to the same transcription \(s\). Let the interpretation values
be distinct \(i_1,i_2\), and put \(I(e_1)=i_1\),
\(I(e_2)=i_2\). Then \(f(e_1)=f(e_2)\) holds but
\(I(e_1)=I(e_2)\) does not. Equality after information loss
therefore supplies no identity of interpretation without an additional
preservation condition.

\emph{Grammatical form and semantic type.} Take two distinct tokens
\(f,c\), let \(H\) contain \((f,c)\), and let
\(A=\{c\}\). Both \(H(f,c)\) and \(A(c)\) are true,
while \(A(f)\) is false. Even making \(H\) the universal
relation, so that it is an equivalence relation of shared form, leaves
the counterexample intact. Semantic category does not follow from
that relation alone.

\emph{Possession and quantity.} The stipulated sets have
\(H_1\subseteq H_0\), \(|H_0|=10\) and \(|H_1|=9\).
The members of \(H_0\) absent from \(H_1\) form exactly
\(\{10\}\), whose size is one. Consequently, losing the
property of possessing ten items is consistent with losing only one.
The related distinction about eyes has the elementary numerical
counterexample \(|E|=2\): it satisfies \(|E|\ne1\), but
not \(|E|=0\). A denial of \emph{exactly one} is not a
denial of \emph{at least one}.

\emph{Per-item and uniform method.} Let the domain be
\(\{a,b\}\), interpret \(L\) as \(\{a\}\) and
\(D\) as \(\{b\}\). Each object satisfies one disjunct,
so \(\forall x(L(x)\lor D(x))\) is true. The instance at
\(b\) falsifies \(\forall xL(x)\), and the instance at
\(a\) falsifies \(\forall xD(x)\). Hence their disjunction
is false. The valid reverse implication follows because either
universal disjunct supplies, for every object, a true disjunct of
\(L(x)\lor D(x)\). It is the unlicensed forward movement
of the disjunction across the quantifier that the example exposes.
\end{proof}
\end{annexitem}

\begin{annexitem}{annex:p18.rd.independent-cases}{thm:p18-rd-independent-cases}{Countermodels, premise dependence and qualified conclusions}
\begin{proof}
\emph{Ownership and parenthood.} Use three distinct domain elements
\(a,d,b\), with \(O=\{(a,d)\}\) and
\(F=\{(d,b)\}\). The two premises of
Example~\ref{def:p18-rd-ownership} are true but \(F(d,a)\)
is false. The witnesses \(x=d,y=b\) verify the retained
existential conclusion in this model and in every model of those
premises. For the contrasting identity case, if \(d\) and \(e\)
denote the same object then \((d,b)\) and \((e,b)\) are the
same ordered pair. Membership in the extension of \(F\) is
therefore preserved. Sharing a property, unlike identity, need not
preserve every relational predication.

For the stated small-number example, \(9<10\) and
\(9\cdot9=81\geq10\), whereas \(2<10\) and
\(2\cdot2=4<10\). Thus neither closure under multiplication
nor universal failure of smallness follows from the factors' being
small under that one fixed threshold.

\emph{Shared consequence without identity.} In a two-object domain
let \(P=\{a\}\), \(Q=\{b\}\) and \(T=\{a,b\}\).
Both \(\forall x(P(x)\to T(x))\) and
\(\forall x(Q(x)\to T(x))\) hold, but the extensions of
\(P,Q\) are different, and their instances need not be the same
objects. This models the limited information supplied by a shared
characteristic, such as whiteness of two different things.

\emph{Qualifications and relation arguments.} Interpret \(G\)
so that \(G(a,c_1)\) is true and \(G(a,c_2)\) false, where
\(c_1\ne c_2\). This satisfies the two qualified premises
without satisfying \(\neg G(a,c_1)\). The apparent contradiction
arises only after erasing the distinguishing argument. For the
arithmetic relation \(D(x,y)\) defined by \(x=2y\),
\(D(2,1)\) holds because \(2=2\cdot1\), whereas
\(D(2,3)\) fails because \(2\ne6\). Both stated premises
are consequently true together. They cannot be the contradictory
of each other with the comparison fixed.

\emph{Quantified target.} In any interpretation with \(\neg P(a)\),
the denotation of \(a\) witnesses \(\exists x\neg P(x)\)
and falsifies \(\forall xP(x)\). It need not falsify every
\(P\)-instance: a second object \(b\) can satisfy \(P\).
Conversely, in the two-object interpretation \(P=\{a\}\),
\(P(a)\) and \(\exists x\neg P(x)\) are both true.
Thus a favorable instance does not refute an existential counterexample.

\emph{Converse and inverse.} At \(p=0,q=1\), the implication
\(p\to q\) is true. The extra premise \(q\) is true but
\(p\) false, defeating the converse pattern. Also \(\neg p\)
is true but \(\neg q\) false, defeating the inverse pattern.
For valid contraposition, any valuation satisfying \(\neg q\)
has \(q=0\). If it also satisfies \(p\to q\), it cannot
have \(p=1\), because that would falsify the implication. Hence
it satisfies \(\neg p\). If instead a warranted converse
\(q\to p\) is supplied together with \(q\), its truth
forces \(p\). The validity of that repair leaves the warrant
for the added converse as a separate question.

\emph{The original point as its own premise.} Every valuation
satisfying \(p\) satisfies \(p\), so
\(\{p\}\models p\). A valuation with \(p=0\) witnesses
\(\varnothing\not\models p\). Nothing in the first entailment
provides an independent reason to eliminate that valuation from the
assessment of an unsupported \(p\). Whether an argument has
independent warrant depends on the stated premise support and its
dialectical role, not merely on this truth-preserving sequent.

\emph{Irrelevant addition to an inconsistent background.} No valuation
satisfies \(\{p,\neg p\}\); adjoining \(h\) leaves that
fact unchanged. Every valuation of that background satisfies \(h\)
and \(\neg h\) vacuously, so the procedure cannot single out
\(h\) for rejection. More generally, suppose \(\Delta\)
is satisfiable but \(\Delta\cup\{h\}\) is not. A valuation
satisfying \(\Delta\) cannot make \(h\) true, or it would
satisfy the enlarged set. Thus every such valuation makes \(h\)
false and \(\Delta\models\neg h\). Satisfiability of
\(\Delta\) ensures this conclusion is nonvacuous. It neither
establishes that the actual world satisfies \(\Delta\) nor
supplies a causal interpretation of the incompatibility.

\emph{One negative answer to two components.} Put \(P(a)=1\)
and \(P(b)=0\). Then \(\neg(P(a)\land P(b))\) is true,
but \(\neg P(a)\land\neg P(b)\) false. For any two
Boolean values, failure of their conjunction means at least one is
zero; hence at least one corresponding negation is true. Conversely
either false component makes the conjunction false. This proves
\[
 \neg(P(a)\land P(b))\quad\Longleftrightarrow\quad
 \neg P(a)\lor\neg P(b),
\]
and explains exactly which conclusion the negative answer supports.
\end{proof}
\end{annexitem}
